% What's the problem?
% Why is it a problem? Research gap left by other approaches?
% Why is it important? Why care?
% What's the approach? How to solve the problem?
% What's the findings? How was it evaluated, what are the results, limitations, 
% what remains to be done?

% XXX Summary
\emph{Summary:}
\dots

% Motivation and intended learning outcomes
\emph{Avsedda lärandemål:}
Efter att ha genomfört detta kapitel ska du kunna:
\begin{itemize}
  \item Förklara varför funktioner gör program kortare, lättare att läsa och
    lättare att ändra. % Lo01
  \item Skriva egna funktioner med parametrar och returvärden och anropa
    dem. % Lo03
  \item Skilja på parametrar, argument och returvärden samt på lokala
    variabler och variabler i huvudprogrammet. % Lo03
  \item Dela upp ett program i funktioner så att kodupprepning undviks. % Lo01
  \item Dokumentera funktioner med docstrings enligt PEP~257. % Lo07
\end{itemize}

% XXX Prerequisites
\emph{Prerequisites:}
\dots

% XXX Reading material
\emph{Reading:}
\dots
